\chapter{Dataset Preparation and Model Development}
\label{ch:data-models}

This chapter describes the evidence pipeline that converts heterogeneous
third-party imagery into a controlled detector comparison. It first separates
the historical project assets from the fire-only research dataset, then
documents annotation and leakage controls, frozen data partitions, training,
model selection, export and conversion-fidelity checks. Predictive results are
reported in Chapter~\ref{ch:results}; the purpose here is to make clear which
inputs, decisions and artifacts produced those results.

\section{Dataset provenance and historical development}
\label{sec:data-models:1}

\subsection{Historical task definition}

The project did not begin with the final one-class experiment. The historical
Sentinel and Scout deployments used models and datasets that included fire and
smoke-related classes. Those artifacts are retained as part of the engineering
history, but they are not treated as the source of the controlled six-model
comparison. Live smoke was removed from the formal validation scope because a
safe and repeatable physical smoke source was not available. The controlled
task therefore uses one localized target class, \emph{fire}, while smoke-only
imagery is excluded from the positive class.

The recovered base export identifies itself as Drone-Fire-Smoke version~1. Its
local metadata records 41,880 images, 640-pixel stretch preprocessing and
flip, brightness and exposure augmentation. The public dataset is distributed
through Roboflow Universe under CC~BY~4.0 \cite{DroneFireSmokeDataset}. I did
not collect these images, and the available export does not provide reliable
original video or capture-session identifiers. Consequently, the image count
cannot be interpreted as the number of independent scenes.

Historical development also introduced environmental negative images to expose
the detector to fire-like visual confounders. Filename and content mapping
recovered images from the Weather Image Recognition dataset, which contains
multiple weather categories and declares a CC0 licence
\cite{BhathenaWeatherDataset}, and from the Weather Classification dataset,
which also declares CC0 \cite{GuptaWeatherClassification}. The former records
the Weather Phenomenon Database as its source \cite{Xiao2021WEAPD}. A further
711 images were mapped to the Stanford Background Dataset, while 21 historical
files had no surviving source match. The Stanford and unresolved groups were
excluded from the research freeze because their reuse conditions or provenance
were not sufficiently established. This exclusion is a licensing and
traceability decision, rather than an assertion that their image content was
invalid.

\subsection{Fire-only curation decision}

Label inspection divided the 41,880-image base export into 20,584 images with
at least one fire annotation, 21,190 smoke-only or other images, and 106 empty
or background images. Fire-positive images retained only class-0 fire boxes.
Smoke-only images were excluded rather than silently relabelled as negative,
because their source annotations were not designed to establish the absence of
visible fire. The approved environmental sources contributed 2,583 candidate
background images. Exact hashing removed 53 duplicate negative copies, leaving
2,530; with the 106 base backgrounds, the provisional dataset contained
20,584 positive and 2,636 background images.

The distinction between a negative example and erroneous data is important.
Fog, cloud, glare, sunrise, sandstorm and similar scenes are legitimate
fire-free observations. Retaining them permits image-level false-alert
behaviour to be measured. Their inclusion does not prove that these conditions
represent the future industrial domain, and the positive and negative images
come from different source collections, so source-specific cues may remain.

\subsection{From provisional data to the frozen research manifest}

The provisional 23,220-image manifest passed structural checks, but further
analysis showed that three inferred source groups contained 16,364 records,
or 70.5\% of the provisional collection. Allowing those groups to dominate
separate partitions would create a strong source/domain confound. The larger
manifest was therefore preserved as an intermediate artifact and was not used
for the final comparison.

Figure~\ref{fig:dataset-lineage} summarizes the controlled lineage. I first
resolved ten flagged annotations into immutable source-family manifest version
10, which contained 23,214 images and 29,095 fire boxes. I then excluded the
three dominant mega-groups, retained 6,850 records from 2,628 inferred groups,
and deterministically capped every retained group at 50 records using seed
\texttt{20260729}. This produced a 5,185-record candidate with more balanced
positive and background support. Whole inferred groups, rather than individual
images, were assigned to train, validation and test targets.

\begin{figure}[H]
  \centering
  \resizebox{0.94\textwidth}{!}{\begin{tikzpicture}[
  node distance=7mm,
  box/.style={draw,rounded corners=2pt,align=left,text width=0.76\textwidth,
    minimum height=10mm,inner sep=3mm,fill=blue!4},
  decision/.style={box,fill=orange!9},
  tag/.style={draw,rounded corners=2pt,align=center,text width=18mm,
    minimum height=10mm,font=\scriptsize\bfseries,fill=gray!10},
  arrow/.style={-{Latex[length=2.2mm]},thick}
]
  \node[box] (source) {\textbf{Historical source pool.} 41,880 augmented fire/smoke images, supplemented by recovered environmental negative-image folders.};
  \node[tag,left=7mm of source] {SOURCE};
  \node[decision,below=of source] (policy) {\textbf{Eligibility policy.} Retain images with at least one fire box; exclude smoke-only images; admit licensed fire-free hard negatives.};
  \node[tag,left=7mm of policy] {POLICY};
  \node[box,below=of policy] (prov) {\textbf{Provisional manifest.} 23,220 images: 20,584 positive images and 2,636 background images.};
  \node[tag,left=7mm of prov] {MANIFEST};
  \node[decision,below=of prov] (audit) {\textbf{Integrity audit.} Reconstruct source families, adjudicate annotations, screen exact and perceptual duplicates, and control splits at group level.};
  \node[tag,left=7mm of audit] {AUDIT};
  \node[box,below=of audit] (v10) {\textbf{Immutable source-family version 10.} 23,214 images and 29,095 boxes.};
  \node[tag,left=7mm of v10] {BASELINE};
  \node[decision,below=of v10] (curate) {\textbf{Controlled curation.} Exclude three dominant mega-groups, cap retained groups at 50, remove ambiguous records, and quarantine five unreliable source families.};
  \node[tag,left=7mm of curate] {CURATION};
  \node[box,below=of curate] (final) {\textbf{Frozen research manifest version 7.} 4,991 images, 3,400 fire boxes and 2,615 inferred groups; group-preserving split of 3,558 train, 700 validation and 733 one-time test images.};
  \node[tag,left=7mm of final] {FROZEN};

  \draw[arrow] (source) -- (policy);
  \draw[arrow] (policy) -- (prov);
  \draw[arrow] (prov) -- (audit);
  \draw[arrow] (audit) -- (v10);
  \draw[arrow] (v10) -- (curate);
  \draw[arrow] (curate) -- (final);
\end{tikzpicture}
}
  \caption{Dataset lineage from the historical fire/smoke export to the frozen
  fire-only research split. Every transformation created a new manifest; the
  source images, source annotations and previous manifests were preserved.}
  \label{fig:dataset-lineage}
\end{figure}

A split- and box-size-stratified review then led to the exclusion of six
ambiguous positive records and the removal of 28 non-fire or redundant boxes
from nine records. A new sample exposed systematic errors in augmented sibling
images, so five source families were quarantined as complete units, removing
188 further records. The resulting version-7 freeze contains 4,991 records,
3,400 retained fire boxes and 2,615 inferred source groups. This 78.5\%
reduction relative to version 10 is an independence and quality control; it
does not mean that 78.5\% of the source imagery was intrinsically false.

Table~\ref{tab:data-provenance} records the sources that remain in the freeze.
The URLs and recorded licence strings are stored per record. The original
dataset pages remain the attribution authority; the manifest is the local
traceability record.

\begin{table}[H]
  \centering
  \caption{Recorded provenance of the frozen version-7 dataset. Counts refer
  to retained records after grouping, capping and adjudication.}
  \label{tab:data-provenance}
  \small
  \begin{tabularx}{\textwidth}{@{}XrX@{}}
    \toprule
    Recorded source & Records & Recorded reuse status \\
    \midrule
    Drone-Fire-Smoke version 1 & 2,461 & CC BY 4.0; attribution required \\
    Weather Image Recognition / WEAPD derivative & 1,366 & CC0 as recorded on the Kaggle page \\
    Weather Classification & 1,164 & CC0 as recorded on the Kaggle page \\
    \midrule
    Total & 4,991 & Third-party imagery; not collected for this thesis \\
    \bottomrule
  \end{tabularx}
\end{table}

\section{Annotation policy, data health and leakage control}
\label{sec:data-models:2}

\subsection{Class and annotation representation}

The controlled experiment has one target class, \emph{fire}. It is represented
as class index 0 in YOLO text labels and category identifier 1 in COCO and
TFRecord exports. A positive record must contain at least one retained fire
box; a background record must contain none. The same manifest was materialized
into framework-specific representations so that the candidate frameworks did
not receive different semantic partitions.

For each YOLO row with normalized centre $(x_c,y_c)$ and normalized width and
height $(w,h)$, the audit reconstructed the edges
\begin{equation}
  x_{\min}=x_c-\frac{w}{2},\quad
  x_{\max}=x_c+\frac{w}{2},\quad
  y_{\min}=y_c-\frac{h}{2},\quad
  y_{\max}=y_c+\frac{h}{2}.
  \label{eq:yolo-box-edges}
\end{equation}
Valid rows require finite values, class index zero, $x_c,y_c\in[0,1]$,
$w,h\in(0,1]$, and all reconstructed edges within the image to a numerical
tolerance of $10^{-9}$. The audit also checked record/image/label existence,
positive image dimensions, duplicate rows and consistency between record type
and retained box count.

The earlier 23,220-record audit found no malformed row or geometric error, but
the deterministic 45-image visual queue identified 35 acceptable, four
questionable and six apparently incorrect cases. Those proportions are not an
estimate of population label error: the queue was deliberately stratified by
split and box size, and one reviewer first judged contact-sheet renders.
Full-resolution adjudication removed six records and 16 boxes when producing
version 10. The subsequent curated lineage repeated semantic review, removed
the further ambiguous records and boxes described above, and quarantined the
five affected source families. All 50 records in the final version-7
split/size-stratified queue were accepted. This final check still cannot prove
that all remaining boxes are semantically correct.

\subsection{Image and group integrity}

The final audit decoded all 4,991 referenced images. It reported no corrupt or
unreadable file, no defined annotation-format problem, no invalid class or box
geometry, and no inferred group identifier spanning more than one split.
Seventeen exact-duplicate sets remain, but every set is contained within a
single split. A difference-hash screen retained 456 cross-split candidates
because low-detail sky, fog, glare and flame images can collide under a compact
perceptual hash. A stabilized 108-pair, high-risk manual queue contained no
remaining visually same-scene pair.

Keeping inferred groups intact is intended to reduce leakage from augmented
or neighbouring observations. Closely related samples across training and
evaluation partitions can inflate apparent generalization
\cite{Tampu2022Leakage}. The cited study concerns a different imaging domain;
it supports the dependence mechanism rather than providing a numerical
baseline here. Because the source export lacks authoritative session IDs,
filename parsing, hashes and manual review may over-group unrelated images or
under-group related ones. The split is therefore described as
\emph{leakage-controlled}, not as proven capture-session independent.

\subsection{Object-size boundary}

The training split contains 2,505 fire boxes. The validation split contains
392 boxes and the one-time test split contains 503. At stored image resolution,
none of the validation boxes is below the COCO $32^2$-pixel small-object
boundary, while the test split contains only three such boxes. Standard
COCO-small average precision is therefore unavailable for validation and too
sparsely supported for a general test-set conclusion \cite{Lin2014COCO}.
Descriptive recall is additionally stratified by the ground-truth box area as
a fraction of image area: below 1\%, 1--5\%, 5--20\% and at least 20\%.
These relative bins do not recreate the image-space size of the small lighter
flame used in the stationary laboratory demonstration.

\section{Frozen research splits and external data}
\label{sec:data-models:3}

\subsection{Partition roles and reproducibility anchors}

The final split was assigned at inferred-group level toward a deterministic
70/15/15 target. Table~\ref{tab:dataset-splits} gives the resulting composition.
The training partition fits model parameters. The validation partition selects
checkpoints, confidence thresholds and the validation interpretation. The test
partition was sealed until those choices and their hashes were frozen.

\begin{table}[H]
  \centering
  \caption{Frozen fire-only version-7 composition. A background image contains
  no retained fire box.}
  \label{tab:dataset-splits}
  \begin{tabular}{lrrr}
    \toprule
    Split & Positive images & Background images & Total images \\
    \midrule
    Training & 1,731 & 1,827 & 3,558 \\
    Validation & 319 & 381 & 700 \\
    Test & 354 & 379 & 733 \\
    \midrule
    Total & 2,404 & 2,587 & 4,991 \\
    \bottomrule
  \end{tabular}
\end{table}

The authoritative records JSONL has SHA-256
\path{478803efb7f8f4f6c3fe74d58658958bccfa2d0f58ec50a1a5c16535e393a7dd}.
The COCO annotations are separately hashed for training, validation and test.
The validation and test annotation hashes are respectively
\path{1a93f800d63bf355bb586039da245cd612cbdd7052102cbfd6da1409312ab7c6}
and
\path{651e2782db96cbbdf2ba718494f7b57447f96aea6d7719fac350132f1b72a06d}.
These anchors allow generated YOLO, COCO and TFRecord materializations to be
checked without treating generated copies as new datasets.

The one-time test split was opened on 9 August 2026 only after the six
validation outputs, checkpoint identities, confidence thresholds, evaluator
hash and selection interpretation were frozen. All six candidates were then
evaluated once at their validation-selected thresholds. No test-driven
threshold tuning, checkpoint reselection or retraining followed. Chapter~7
reports the predictive outcomes and the rank reversal observed between
validation and test; the present chapter records only the procedural boundary.

\subsection{External-evaluation eligibility}

External UAV and industrial-like sources are reserved for RQ5. They did not
enter training, validation or the one-time internal test. At the final evidence
freeze, no external inference result was admitted. FASDD-UAV has a
publisher-reported CC~BY-SA~4.0 licence and 25,097 UAV samples, but the exact
archive was not acquired, hashed, decoded, audited for overlap and
frozen locally. These publisher values are not substituted for verified local
counts.

The MmodalFire archive is retained and hash-verified, and all 28 videos were
decoded to their ends without a decoder failure. However, only eight of 28
video ranges passed the strict consolidated timestamp-alignment rule; the
remaining ranges did not complete event/condition review and overlap checks.
Optical
character recognition is used only as review assistance. Until eligible
non-overlapping manifests are frozen, these sources cannot contribute model
metrics. Any later outcomes should be described as external UAV-domain and
industrial-domain proxy evaluations, not as operational industrial-site
validation.

\section{Candidate detectors and training protocol}
\label{sec:data-models:4}

\subsection{Controlled candidate matrix}

The experiment compares six pretrained single-stage detectors. Four YOLOv8
configurations form a $2\times2$ capacity/resolution comparison: nano and small
variants at 416 and 640 pixels. SSD MobileNet V2 FPNLite and
EfficientDet-Lite0 provide mobile-oriented non-YOLO comparators
\cite{UltralyticsYOLOv8,Liu2016SSD,Sandler2018MobileNetV2,Tan2020EfficientDet}.
The screening seed is \texttt{20260729} for every candidate. Input resolution
is part of the candidate configuration; comparisons between 416, 640 and
320 pixels are therefore deployment-oriented comparisons rather than pure
architecture ablations.

\begin{table}[H]
  \centering
  \caption{Six controlled detector configurations and their recorded training
  completion. Predictive and runtime results are presented in Chapter~7.}
  \label{tab:candidate-training}
  \scriptsize
  \begin{tabularx}{\textwidth}{@{}p{0.13\textwidth}p{0.21\textwidth}p{0.11\textwidth}p{0.10\textwidth}p{0.12\textwidth}X@{}}
    \toprule
    ID & Model / framework & Input & Batch & Completion & Training lineage \\
    \midrule
    N416 & YOLOv8n / PyTorch & 416 px & 16 & 48 epochs & Early stopping under the 100-epoch, patience-20 policy \\
    N640 & YOLOv8n / PyTorch & 640 px & 4 & 87 epochs & Optimizer-state resume from epoch 51 \\
    S416 & YOLOv8s / PyTorch & 416 px & 8 & 33 epochs & Early stopping under the same policy \\
    S640 & YOLOv8s / PyTorch & 640 px & 4 & 33 epochs & Early stopping under the same policy \\
    SSD320 & SSD MobileNet V2 FPNLite / TensorFlow OD API & 320 px & 8 & 44,500 steps & One-class pipeline; corrected bounded validation signature \\
    EDL0 & EfficientDet-Lite0 / TFLite Model Maker & 320 px & 4 & 100 epochs & Full-state recovery through checkpoints 50 and 90 \\
    \bottomrule
  \end{tabularx}
\end{table}

All candidates use the same version-7 record allocation. Pretrained
initialization, environment locks, dataset hashes, checkpoints and failure
lineage are retained per run. The four YOLO configurations ran in an isolated
Windows Python~3.12 environment with PyTorch~2.12, Ultralytics~8.4.45 and an
NVIDIA RTX~3050 Laptop GPU. Their starting weights were hash-verified. The
training request specified 100 maximum epochs, early-stopping patience 20,
automatic mixed precision and the largest stable recorded batch for each
configuration.

SSD training used TensorFlow~2.15.1 and a pinned TensorFlow Object Detection
API revision under WSL Python~3.11. The pipeline used a 320-pixel input,
batch size 8, 445 steps per epoch, 100 maximum epochs and a cosine-decay
momentum schedule. An initial common validation attempt created excessive
TensorFlow retracing as image shapes changed, reaching 7.5~GB WSL resident
memory and severe swap activity. That attempt was preserved. The predictor was
corrected with one dynamic-spatial-dimension signature; the bounded retry used
two concrete traces and completed normally. This was a runtime/evaluator
correction, not a change to the model or ground truth.

EfficientDet-Lite0 used TFLite Model Maker~0.4.3 and TensorFlow~2.8.1 in a
separate WSL Python~3.9 environment. The legacy loader initially failed on
genuine zero-box background examples. Instead of deleting 2,587 background
records or generating artificial boxes, I introduced a source-controlled
empty-safe mask with the same semantics for non-empty examples. A one-epoch
gate confirmed non-zero training and validation box loss before the full run.
The full-state lineage, including optimizer iterations, was preserved through
interruption and resume; checkpoint 90 supplied the final ten epochs, and the
epoch-100 run ended normally.

The six runs constitute a single-seed screen. A three-seed finalist stage was
originally planned, but it was not executed before the selection freeze.
Consequently, between-seed variance cannot be estimated and no claim of
statistical superiority is based on the training runs. Failed, interrupted and
superseded attempts remain part of the engineering record and are excluded
from accepted predictive comparisons unless a protocol explicitly adopts a
recovered run.

\section{Validation-based selection and held-out evaluation}
\label{sec:data-models:5}

Each native framework generated COCO-format predictions down to confidence
0.001 on the same 700 validation images. A common evaluator computed COCO AP
over IoU thresholds 0.50:0.05:0.95 and AP50. For each candidate, the reporting
confidence threshold was selected from 0.01 to 0.99 by maximizing object-level
F1 at IoU 0.50. A numerical tie was resolved by higher recall and then by the
lower threshold. The selected threshold was frozen before the test split was
opened.

For a confidence threshold $t$, greedy class-aware matching at IoU 0.50 gives
true positives $TP(t)$, false positives $FP(t)$ and false negatives $FN(t)$.
The selection quantities are
\begin{align}
  P(t) &= \frac{TP(t)}{TP(t)+FP(t)}, \\
  R(t) &= \frac{TP(t)}{TP(t)+FN(t)}, \\
  F_1(t) &= \frac{2P(t)R(t)}{P(t)+R(t)}.
  \label{eq:threshold-f1}
\end{align}
An object detector has no unique count of true-negative boxes. Negative
behaviour is therefore reported as the fraction of background images with at
least one prediction at $t$, together with false-positive counts, rather than
as a fabricated box-level specificity.

Model selection considered AP50--95, validation-selected F1, recall,
precision, background rejection, observed host latency, artifact size and
demonstrated device feasibility as separate dimensions. No post-hoc weighted
score was introduced. A candidate is Pareto-efficient when another candidate
cannot match or improve every selected dimension and strictly improve at least
one. Sentinel interpretation emphasizes recall subject to its device
constraints; Scout interpretation emphasizes precision and rejection of
background alerts. These role preferences do not convert laptop screening
into evidence of deployment.

The selection freeze was written before the test opening and records all six
thresholds and checkpoint lineages. The test runner verified the freeze and
test-annotation hashes, refused silent overwrites and retained per-image
predictions and timing. A Windows-to-WSL path-normalization defect stopped the
first invocation after the first model's inference. The raw predictions were
preserved, the artifact-level path handling was corrected, and evaluation
continued without rerunning or changing those predictions. This recovery is
part of the evidence lineage rather than an unreported clean run.

\section{Export, optimization and deployment fidelity}
\label{sec:data-models:6}

\subsection{Fidelity before performance}

A converted model is not assumed to be equivalent to its native checkpoint.
The deployment registry requires comparison on the same 64 hashed validation
calibration images before a converted artifact can enter device timing. Native
and converted predictions use declared preprocessing and decoding. Within each
image and class, detections are matched at box IoU at least 0.50. The frozen
conversion rule requires a detection match rate of at least 0.95 and a mean
absolute matched-confidence difference no greater than 0.02. A failure on
either condition blocks performance ranking for that converted artifact.

The four YOLO checkpoints were exported as fixed batch-1, fixed-input FP32 ONNX
graphs at opset 12 without simplification. ONNX~1.22 accepted all graph
structures, but graph validity is distinct from prediction fidelity. Table
\ref{tab:export-fidelity} records the current prepared form and disposition.
The detailed counts belong to Chapter~7; here, \texttt{FIDELITY\_FAILED} means
that the frozen conversion gate was not met.

\begin{table}[H]
  \centering
  \caption{Prepared candidate artifacts and conversion-fidelity disposition.
  These are laptop preparation results, not edge-device benchmarks.}
  \label{tab:export-fidelity}
  \scriptsize
  \begin{tabularx}{\textwidth}{@{}p{0.20\textwidth}p{0.18\textwidth}p{0.15\textwidth}X@{}}
    \toprule
    Candidate & Prepared form & Fidelity disposition & Permitted next step \\
    \midrule
    YOLOv8n-416 & FP32 ONNX, opset 12 & Pass & TensorRT build/validation on Sentinel; ONNX Runtime attempt on Scout \\
    YOLOv8n-640 & FP32 ONNX, opset 12 & Failed & Independent TensorRT candidate may be built; current ONNX is rejected \\
    YOLOv8s-416 & FP32 ONNX, opset 12 & Pass & TensorRT build/validation on Sentinel; ONNX Runtime attempt on Scout \\
    YOLOv8s-640 & FP32 ONNX, opset 12 & Failed & Independent TensorRT candidate may be built; current ONNX is rejected \\
    SSD MobileNet V2 FPNLite-320 & Float32 TFLite & Failed & Audit another native/conversion path or classify as unsupported \\
    EfficientDet-Lite0-320 & Native TFLite & Not applicable & Attempt only where a compatible TFLite runtime is available \\
    \bottomrule
  \end{tabularx}
\end{table}

The SSD conversion produced the expected number of outputs after its tensor
roles and resize semantics were corrected, but only 87.14\% of detections
matched the native predictions, below the fixed 95\% requirement. Its small
matched-confidence error did not override the unmatched detections. The
EfficientDet artifact is the native TFLite output of its training toolchain, so
a separate source-to-converted comparison is not applicable; device runtime
support and resource cost still require measurement.

\subsection{Historical deployed models versus screened candidates}

The six frozen checkpoints are research-screening artifacts. They must not be
silently substituted for the models already exercised in the two-device
system. The final Sentinel workspace contains a distinct YOLOv8n TensorRT
engine with input shape $[1,3,640,640]$, output shape $[1,6,8400]$ and SHA-256
\path{16f9018ac299c6f525610a877d798392ec26d53fccb84b589ff6533d76589bf2}.
Earlier 416-pixel measurements refer to a different engine with output shape
$[1,6,3549]$ and cannot be applied to this deployment.

The historical Scout deployment uses a distinct TFLite artifact with input
$[1,640,640,3]$, output $[1,7,8400]$ and SHA-256
\path{d533e320081be17ac91fb77c20a9de09f4bffada0d7f88c88326d57cd893bc35}.
The tensor interface is float32 and the seven output channels are consistent
with four box values and three class scores. It belongs to the historical
multi-class development line, rather than to the one-class six-model freeze.

The strengthened exact-device campaign for the six new candidates was not
executed before the final evidence freeze. Its protocol requires 20 warm-up and
200 measured inferences over the same 64-image calibration manifest, one-second
resource sampling, artifact and image hash verification, and separate
preprocessing, runtime, postprocessing and complete-pipeline timing. None of
the twelve candidate/platform pairs has a frozen outcome. The existing
historical Sentinel and Scout measurements establish the demonstrated pair;
they do not fill the six-model device matrix.

\section{Chapter evidence boundary}

This chapter establishes a reproducible data and model-development lineage:
the version-7 dataset is frozen, all six single-seed candidates completed,
validation-based thresholds were frozen, and the test split was evaluated once
without post-test tuning. The chapter does not claim that the dataset is
capture-session independent, that all annotations are error free, that any
model family is universally superior, or that the six candidates have been
validated on both edge devices. External transfer and the strengthened
exact-device matrix remain explicit evidence gaps. These boundaries determine
which numerical results can be interpreted in Chapters~7 and 8.
